% ==============================================================================
% ПАРТИЯ B03-014: DISS-005-B0587 .. DISS-005-B0656
% Глава 1: Блокировка работы БАВР по РНМ при внешних двухфазных замыканиях
% Замечания: DISS-005-C0063 .. DISS-005-C0065
% Иллюстрации: image27.png .. image40.png
% ==============================================================================

% DISS-005-B0587
\subsection{Блокировка работы БАВР по РНМ при внешних двухфазных замыканиях}
\label{sec:ch1_external_two_phase_blocking}

% DISS-005-B0588
Иногда сложности функционирования определяются формой характеристики срабатывания РНМ. Так, например при двухфазных замыканиях за трансформатором -- вектор может оказываться на границе зоны срабатывания, что приводит к блокировке. В следствии чего произошло отключение двигательной нагрузки по минимальному напряжению. Далее приведён разбор поведения БАВР в таком случае, описание причин данного функционирования и рекомендации по корректировке параметров для улучшения функционирования.

% DISS-005-B0590
\subsubsection{Функционирование БАВР}

% DISS-005-B0591
В данной сети произошло аварийное событие, продолжительность начального этапа составляла порядка 300-350~мс (от 1~с до 1,35~с), что видно по напряжению на 1СШ (рис. 1а). Здесь напряжение восстановилось практически до доаварийных значений. Далее аварийные процессы не оказывали влияния на данную секцию, вероятно, ввиду отключения внешних связей между секциями.

% DISS-005-B0592
При этом, на второй секции, где находилась двигательная нагрузка, после завершения этапа 1, началось восстановление нормально режима -- этап 2 (от 1,35~с до 2,05~с). Ввиду наличия мощной двигательной нагрузки (6,3~МВт) данный процесс занимает продолжительное время, и он не успел завершиться, так как произошло повторное замыкание во внешней сети по другим фазам -- этап 3 (от 2,05~с до 2,5~с). Ввиду того, что длительное время напряжение на секциях шин находилось ниже допустимого отработала защита по минимальному напряжению и отключила двигатель.

% DISS-005-B0594 .. B0604 (DISS-005-FIG-024, 025, 026 / DISS-005-TAB-018)
\begin{figure}[htbp]
    \centering
    \begin{minipage}{0.9\textwidth}
        \centering
        \includegraphics[width=\textwidth]{image27.png}
        \centerline{\small а)}
    \end{minipage}\\[1em]
    \begin{minipage}{0.9\textwidth}
        \centering
        \includegraphics[width=\textwidth]{image28.png}
        \centerline{\small б)}
    \end{minipage}\\[1em]
    \begin{minipage}{0.9\textwidth}
        \centering
        \includegraphics[width=\textwidth]{image29.png}
        \centerline{\small в)}
    \end{minipage}
    \caption[Рисунок 1. Полный вид осциллограммы произошедшего события]{Рисунок 1. Полный \commentnote{DISS-005-C0063}{Евдаков Алексей}{t00051\_240207\_2327\_\_0021}{\commentnote{DISS-005-C0064}{Евдаков Алексей}{«ДКС-1 ЗРУ-10кВ СШ1-3» в 23:27 (07.02.2024)}{вид }}осциллограммы произошедшего события а) 1СШ б) 3СШ*, в) указание разделения этапов.\newline{\small Данные файлы находятся в папке «\_вспомогательные файлы\textbackslash{}ДКС-1\textbackslash{}Блокировка работы БАВР по РНМ при внешних двухфазных замыканиях»}}
    \label{fig:ch1_img27_29}
\end{figure}

% DISS-005-B0605
*Примечание: БАВР предусматривает работу по двум секциям, поэтому данная секция далее везде именуется как «2СШ».

% DISS-005-B0606
Просадка линейных напряжений по 2СШ была значительно сильнее, и первым возникли условия для пуска именно по этой секции -- защита 27/27s (рис. 2). Время реакции органов БАВР составили 3,125~мс. Однако ввиду особенностей протекания тока (нахождения РНМ у границы), не возникло разрешения от реле направления мощности. Затем напряжение и по первой секции шин опустилось ниже порога срабатывания, и просадка считалась синхронной ввиду чего работа БАВР была заблокирована (пропали сигналы готовности БАВР).

% DISS-005-B0607 .. DISS-005-B0608 (DISS-005-FIG-027 / DISS-005-TAB-019)
\begin{figure}[htbp]
    \centering
    \includegraphics[width=0.85\textwidth]{image30.png}
    \caption{Рисунок 2. Функционирование БАВР, сигналы пуска и блокировки}
    \label{fig:ch1_img30}
\end{figure}

% DISS-005-B0609
Процесс отсутствия готовности работы БАВР длился порядка 1,45~с (рис. 3), из них 0,35~с длилась сама авария (блокировка БАВР по синхронной просадке), и потом ещё 1,1~с задержка на готовности БАВР (данный параметр является уставкой задаваемой в вкладке «Параметры БАВР»).

% DISS-005-B0610 .. DISS-005-B0611 (DISS-005-FIG-028 / DISS-005-TAB-020)
\begin{figure}[htbp]
    \centering
    \includegraphics[width=0.85\textwidth]{image31.png}
    \caption{Рисунок 3. Сигналы готовности БАВР}
    \label{fig:ch1_img31}
\end{figure}

% DISS-005-B0612
Отметим, что продолжительность этапа 1 (аварии) и последующего восстановление готовности БАВР примерно совпало с уставкам РЗА, установленными в данной сети. Данные процессы не связаны между собой, и являются лишь совпадением.

% DISS-005-B0613
\subsubsection{Тип аварии}

% DISS-005-B0614 .. DISS-005-B0616
Проведена оценка типа аварии (выявления фаз), в которых произошло внешнее замыкание с помощью программного обеспечения «Проект "РЗА"»:

\url{https://pro-rza.ru/models/ensys/}

Примечание: Отметим, что не осуществляется точное моделирование, оценивается лишь характер векторов.

% DISS-005-B0617
Как можем видеть из векторной диаграммы на основе осциллограммы, на этапе 1 присутствует лишь 2 тока фаз А и С (рис. 4а), данный режим можно повторить путём моделирования двухфазного КЗ фаз «В» и «С» в сети 110~кВ (рис. 4б). Данный режим является важным для дальнейших исследований, так как в данных сетях он может приводить к блокировке работы БАВР по РНМ прямой последовательности ввиду особенностей направления протекания токов.

% DISS-005-B0618
На этапе 3, ток фазы С направлен противоположно токам фаз В и С, которые приблизительно равны (рис. 5а), данный режим не удалось в точности восстановить по взаимному направлению векторов тока и напряжения. По предварительной оценке, это однофазное замыкание фазы «В» в сети 110~кВ (рис. 5б).

% DISS-005-B0619 .. DISS-005-B0623 (DISS-005-FIG-029, FIG-030 / DISS-005-TAB-021)
\begin{figure}[htbp]
    \centering
    \begin{minipage}{0.48\textwidth}
        \centering
        \includegraphics[width=\textwidth]{image32.png}
        \centerline{\small а)}
    \end{minipage}\hfill
    \begin{minipage}{0.48\textwidth}
        \centering
        \includegraphics[width=\textwidth]{image33.png}
        \centerline{\small б)}
    \end{minipage}
    \caption{Рисунок 4. Вектора токов и напряжений на этапе 1: а) реальные, б) модельные}
    \label{fig:ch1_img32_33}
\end{figure}

% DISS-005-B0625 .. DISS-005-B0629 (DISS-005-FIG-031, FIG-032 / DISS-005-TAB-022)
\begin{figure}[htbp]
    \centering
    \begin{minipage}{0.48\textwidth}
        \centering
        \includegraphics[width=\textwidth]{image34.png}
        \centerline{\small а)}
    \end{minipage}\hfill
    \begin{minipage}{0.48\textwidth}
        \centering
        \includegraphics[width=\textwidth]{image35.png}
        \centerline{\small б)}
    \end{minipage}
    \caption{Рисунок 5. Вектора токов и напряжений на этапе 3: а) реальные, б) модельные}
    \label{fig:ch1_img34_35}
\end{figure}

% DISS-005-B0631
\subsubsection{Поведение РНМ}

% DISS-005-B0632
На рисунке 6 представлена векторная диаграмма тока и напряжения прямой последовательности на СШ2 на этапе 1 в момент времени 1,2~с (через 200~мс от начала аварии). Как можем видеть, ток отстаёт от напряжения на 30 градусов, что по своему характеру соответствует направлению мощность к нагрузке. В следствии этого, стационарный РНМ воспринимает аварию как внутреннюю и запрещает работу. Функционирования РНМ в данном случае считается корректным.

% DISS-005-B0633 .. DISS-005-B0635 (DISS-005-FIG-033 / DISS-005-TAB-023)
\begin{figure}[htbp]
    \centering
    \includegraphics[width=0.75\textwidth]{image36.png}
    \caption{Рисунок 6. Векторная диаграмма тока и напряжения прямой последовательности на СШ2 на этапе 1}
    \label{fig:ch1_img36}
\end{figure}

% DISS-005-B0637
Безусловно, в переходный момент, наблюдалось колебание вектора РНМ (рис. 7), однако максимальный угол опережения током напряжения составлял лишь 26,6$^\circ$. Стандартная граница срабатывания РНМ в эту сторону по углу составляет 45$^\circ$ с учётом активированной адаптивности по углу максимальной чувствительности в данном случае она была 35$^\circ$, данного значения всё равно недостаточно. Указанный вектор, так же близко проходит к адаптивной границе по амплитуде тока, однако данную уставку коэффициента отстройки равную 0,6~о.е. заменить нельзя, ввиду её отдельной предыдущей настройки для исключения излишнего срабатывания БАВР в режиме пуска двигателя.

% DISS-005-B0638 .. DISS-005-B0640 (DISS-005-FIG-034 / DISS-005-TAB-024)
\begin{figure}[htbp]
    \centering
    \includegraphics[width=0.75\textwidth]{image37.png}
    \caption{Рисунок 7. Траектория движения вектора РНМ относительно зоны срабатывания РНМ на СШ2 на этапе 1}
    \label{fig:ch1_img37}
\end{figure}

% DISS-005-B0641
Ввиду указанных особенностей, для решения проблемы, предлагает использовать адаптивную уставку по углу альфа, которая и предназначена для улучшения работы БАВР при внешних замыканиях при мощной двигательной нагрузке за счёт сужения зоны срабатывания по углу на 10$^\circ$ градусов с стороны «ёмкостной» нагрузки. Таким образом, граница срабатывания составит 35$^\circ$ - 10$^\circ$ = 25$^\circ$, чего на несколько выборок будет достаточно для срабатывания БАВР (рис.~8). Ожидаемое время срабатывания БАВР составит 4,375~мс.

% DISS-005-B0642 .. DISS-005-B0643 (DISS-005-FIG-035 / DISS-005-TAB-025)
\begin{figure}[htbp]
    \centering
    \includegraphics[width=0.85\textwidth]{image38.png}
    \caption{Рисунок 8. Ожидаемое поведение РНМ при активации адаптивной уставки по углу альфа на СШ2 на этапе 1}
    \label{fig:ch1_img38}
\end{figure}

% DISS-005-B0645
\subsubsection{Рекомендации по изменению параметров БАВР}

% DISS-005-B0646
Ниже приведён список рекомендаций по корректировке параметров, который был выдвинут для улучшения функционирования в данных случаях:

% DISS-005-B0647 .. DISS-005-B0648
\begin{itemize}[leftmargin=*,itemsep=0pt]
    \item Увеличить уставку по времени на срабатывание ЗМН используемых для защиты двигателя до 1,7~с или более. Если рассматривать это в более общем случае -- то необходимо согласовывать времена срабатывания РЗА и БАВР, с выдержками задаваемыми для защиты двигателей, так как в данном случае они совпали.
    \item Если это невозможно и в данной сети не бывает режимов с расхождения угла между секциями шин на 13 градусов и более, то можно уменьшить уставку на «Готовность БАВР» до 0,7~с. Данное изменение позволит уменьшить время блокировки БАВР и исключить отключения двигателей по ЗМН. Опять же -- решение проблемы согласования уставок РЗА+БАВР с уставками двигательной нагрузки.
\end{itemize}

% DISS-005-B0649 (DISS-005-FIG-036)
\begin{figure}[htbp]
    \centering
    \includegraphics[width=0.75\textwidth]{image39.png}
    \label{fig:ch1_img39}
\end{figure}

% DISS-005-B0650
\begin{itemize}[leftmargin=*,itemsep=0pt]
    \item Для повышения надёжности функционирования органа РНМ -- активировать адаптивную уставку по углу альфа.
\end{itemize}

% DISS-005-B0651 (DISS-005-FIG-037)
\begin{figure}[htbp]
    \centering
    \includegraphics[width=0.75\textwidth]{image40.png}
    \label{fig:ch1_img40}
\end{figure}

% DISS-005-B0652
Примечание: приведённые скриншоты могу немного отличать в разных версиях ПО.

% DISS-005-B0654
\subsubsection{Примечание по анализу иных осциллограмм}

% DISS-005-B0655
В данный день произошло множество аварийных замыканий во внешней сети. Ряд из них аналогичен приведённому. В ряде из них, нагрузка всё же оставалась в работе, когда не наблюдалось повторного замыкания во внешней сети. А в некоторых указанная уставка по адаптивности не помогла бы, так как максимальный угол РНМ составлял всего \commentnote{DISS-005-C0065}{Евдаков Алексей}{данный режим можно найти на осциллограмме «t00051\_240207\_1414\_\_0090»}{19$^\circ$}. Ещё больше изменять угол зоны считаем нецелесообразным, так как он приблизится к границе и в случае ёмкостного характера, даже с учётом адаптивности, может произойти ложное отпадание РНМ. И, следовательно, при данной реализации проблема не может быть полностью решена изменением уставок.
